\ExerciseSection

\begin{enumerate}
\def\labelenumi{\arabic{enumi})}
\tightlist
\item
  \begin{enumerate}
  \def\labelenumii{\alph{enumii})}
  \tightlist
  \item
    If A is a subset of a topological space X, show that the following statements are equivalent: \(\alpha\) the frontier of A is nowhere dense; \(\beta\) A is the union of an open set and a nowhere dense set; \(\gamma\) A is the difference between a closed set and a nowhere dense set.
  \end{enumerate}
\end{enumerate}

\begin{enumerate}
\def\labelenumi{\alph{enumi})}
\setcounter{enumi}{1}
\tightlist
\item
  Show that if a subset A of X is such that, for each point \(x \in A\) , there is a neighbourhood V of x in X such that \(V \cap A\) is nowhere dense in X, then A is nowhere dense in X.
\end{enumerate}

\begin{enumerate}
\def\labelenumi{\arabic{enumi})}
\setcounter{enumi}{1}
\tightlist
\item
  A subset \(\mathbf{A}\) of a topological space \(\mathbf{X}\) is said to be a thin set if, for each perfect set \(\mathbf{P} \subset \mathbf{X}\) , \(\mathbf{P} \cap \mathbf{A}\) is nowhere dense in \(\mathbf{P}\) .
\end{enumerate}

\begin{enumerate}
\def\labelenumi{\alph{enumi})}
\item
  Every finite union of thin sets is thin.
\item
  The frontier of a thin set is nowhere dense.
\item
  There is a largest perfect set in \(\mathbf{X}\) whose complement is thin.
\end{enumerate}

\begin{enumerate}
\def\labelenumi{\arabic{enumi})}
\setcounter{enumi}{2}
\tightlist
\item
  Let A be a subset of a topological space X, and let D(A) denote the set of all \(x \in X\) such that, for each neighbourhood V of x, the set \(V \cap A\) is not meagre. D(A) is contained in \(\overline{A}\) .
\end{enumerate}

\begin{enumerate}
\def\labelenumi{\alph{enumi})}
\tightlist
\item
  If \(\mathbf{B}\) is a subset of \(\mathbf{X}\) such that \(\mathbf{D}(\mathbf{B}) = \varnothing\) , then
\end{enumerate}

\[
\mathbf {D} (\mathbf {A} \cup \mathbf {B}) = \mathbf {D} (\mathbf {A})
\]

for every subset A of X.

\begin{enumerate}
\def\labelenumi{\alph{enumi})}
\setcounter{enumi}{1}
\item
  D(A) = ∅ if and only if A is meagre. {[}Begin by showing that if (U\_t) is a family of pairwise disjoint open subsets of X and if, for each t, V\_t is a nowhere dense set (relative to X) contained in U\_t, then ∪ V\_t is nowhere dense. Then consider a maximal set M of open subsets of X, each pair of which are disjoint, and such that A ∪ U is meagre for each U ∈ M; the existence of such a maximal set is established by Zorn's lemma. Show that if D(A) = ∅, then A ∩ C W is nowhere dense, where W is the union of the sets of M, and hence show that A is meagre.{]}
\item
  Show that \(\mathbf{D}(\mathbf{A})\) is closed and that \(\mathbf{A} \cap \mathbf{CD}(\mathbf{A})\) is meagre for all \(A \subset X\) {[}show, with the help of b), that \(\mathbf{A} \cap \mathbf{CD}(\mathbf{A})\) is meagre relative to the open subspace \(\mathbf{CD}(\mathbf{A})\) {]}.
\end{enumerate}

\(d\) ) Show that \(\mathbf{D}(\mathbf{A})\) is equal to the closure of its interior. {[}If \(\mathbf{D}'(\mathbf{A})\) is the closure of the interior of \(\mathbf{D}(\mathbf{A})\) , show that \(\mathbf{A} \cap \mathbb{C}\mathbf{D}'(\mathbf{A})\) is meagre, by observing that this set is the union of

\[
\mathbf {A} \cap \mathbb {C D} (\mathbf {A}) \quad \text { and } \quad \mathbf {A} \cap \mathbf {D} (\mathbf {A}) \cap \mathbb {C D} ^ {\prime} (\mathbf {A}). ]
\]

\begin{enumerate}
\def\labelenumi{\arabic{enumi})}
\setcounter{enumi}{3}
\tightlist
\item
  Let \(\mathbf{X}\) and \(\mathbf{Y}\) be two topological spaces and let \(\mathbf{A}\) (resp. B) be a subset of \(\mathbf{X}\) (resp. Y).
\end{enumerate}

\begin{enumerate}
\def\labelenumi{\alph{enumi})}
\item
  \(\mathbf{A} \times \mathbf{B}\) is nowhere dense (resp. perfect) in \(\mathbf{X} \times \mathbf{Y}\) if and only if one of the two sets \(\mathbf{A}, \mathbf{B}\) is nowhere dense (resp. \(\mathbf{A}\) and \(\mathbf{B}\) are closed and one of them is perfect).
\item
  \(\mathbf{A} \times \mathbf{B}\) is thin in \(\mathbf{X} \times \mathbf{Y}\) if and only if \(\mathbf{A}\) and \(\mathbf{B}\) are both thin sets.
\end{enumerate}

\begin{enumerate}
\def\labelenumi{\arabic{enumi})}
\setcounter{enumi}{4}
\tightlist
\item
  \begin{enumerate}
  \def\labelenumii{\alph{enumii})}
  \tightlist
  \item
    Let X and Y be two topological spaces and let A be a nowhere dense subset of \(X \times Y\) . If the topology of Y has a countable base, show that the set of all \(x \in X\) such that the section \(\mathbf{A} \cap (\{x\} \times \mathbf{Y})\) of A at x is not nowhere dense relative to \(\{x\} \times Y\) is a meagre set in X (if U is any non-empty open set in Y, show that the set of all \(x \in X\) such that \(\{x\} \times U\) is contained in the closure of the section of A at x is a nowhere dense set in X).
  \end{enumerate}
\end{enumerate}

\begin{enumerate}
\def\labelenumi{\alph{enumi})}
\setcounter{enumi}{1}
\item
  Show by an example that the result of a) can be false if the topology of Y does not have a countable base (take X to be a Hausdorff space with no isolated points and Y to be the set X with the discrete topology; and take A to be the diagonal of \(X \times Y\) ).
\item
  If B, C are subsets of X, Y respectively and if one of B, C is meagre (in X, Y respectively), then \(B \times C\) is meagre in \(X \times Y\) . Conversely, if \(B \times C\) is meagre in \(X \times Y\) and if the topology of either X or Y has a countable base, then one of the sets B, C is meagre.
\item
  Deduce from c) that if the topology of one of the spaces X, Y has a countable base, then \(\mathrm{D}(\mathbf{B} \times \mathbf{C}) = \mathrm{D}(\mathbf{B}) \times \mathrm{D}(\mathbf{C})\) in the notation of Exercise 3.
\end{enumerate}

\begin{enumerate}
\def\labelenumi{\arabic{enumi})}
\setcounter{enumi}{5}
\tightlist
\item
  A set A in a topological space X is said to be almost open if there is an open set U in X such that both \(U \cap C A\) and \(A \cap C U\) are meagre sets.
\end{enumerate}

\begin{enumerate}
\def\labelenumi{\alph{enumi})}
\item
  The complement of an almost open set is almost open.
\item
  Any countable union of almost open sets is almost open.
\item
  Show that the following statements are equivalent: \(\alpha\) ) A is almost open; \(\beta\) ) there is a meagre subset M of X such that A∩CM is both open and closed in the subspace X---M; \(\gamma\) ) there is a set G⊂A which is a countable intersection of open sets in X, such that A---G is meagre in X; \(\delta\) ) there is a set F⊃A which is a countable union of closed sets in X, such that F---A is meagre in X; \(\zeta\) ) the set D(A)∩D(X---A) (Exercise 3) is nowhere dense in X; \(\theta\) ) the set D(A)∩CA is meagre in X. {[}Use Exercises 3a) and 3b) to show that \(\alpha)\Longrightarrow\zeta)\Longrightarrow\theta)\Longrightarrow\alpha\).{]} d) Let X, Y be two topological spaces and suppose that the topology of either X or Y has a countable base. Then a subset of X × Y of the form \(\mathbf{A} \times \mathbf{B}\) is almost open if and only if either both \(\mathbf{A}\) and \(\mathbf{B}\) are almost open or one of them is meagre {[}use \(\zeta\) ) of \(c\) {]}.
\end{enumerate}

\begin{enumerate}
\def\labelenumi{\arabic{enumi})}
\setcounter{enumi}{6}
\tightlist
\item
  A topological space X is said to be non-meagre if it is not meagre relative to itself.
\end{enumerate}

\begin{enumerate}
\def\labelenumi{\alph{enumi})}
\item
  X is non-meagre if and only if every countable family of dense open sets in X has a non-empty intersection.
\item
  In a non-meagre space, the complement of every meagre set is a non-meagre subspace.
\end{enumerate}

\begin{enumerate}
\def\labelenumi{\arabic{enumi})}
\setcounter{enumi}{7}
\item
  Let X be a topological space in which there is a non-empty open subset A which is a non-meagre subspace. Show that X is non-meagre and that A is not a meagre set relative to X.
\item
  Let X be the subspace of \(R^{2}\) which is the union of \(Q^{2}\) and the line y = 0. Show that X is meagre relative to itself. Deduce that a topological space can have non-meagre subspaces without being non-meagre itself.
\item
  Show that every completely regular pseudo-compact space (§ 1, Exercise 21) is a Baire space (argue as in Theorem 1).
\item
  If X is a Baire space with no isolated points, show that every point of X is a point of condensation of X (Chapter I, § 9, Exercise 17).
\end{enumerate}

¶ 12) A topological space X is said to be totally non-meagre if every non-empty closed subspace of X is non-meagre. A locally compact space is totally non-meagre, and so is a complete metric space.

\begin{enumerate}
\def\labelenumi{\alph{enumi})}
\item
  Show that a totally non-meagre regular space is a Baire space.
\item
  In an accessible totally non-meagre space, a non-empty countable closed set is thin (Exercise 2) and hence has at least one isolated point.
\item
  In a totally non-meagre space X, every subspace A which is the intersection of a countable intersection of open sets in X is totally non-meagre (if F is a subset of A which is closed in A, let \(\overline{F}\) be the closure of F in X, and show that \(\overline{F} \cap C F\) is meagre relative to \(\overline{F}\) ).
\item
  If every point of a topological space X has a neighbourhood which is a totally non-meagre subspace, then X is totally non-meagre.
\end{enumerate}

¶ 13) Let X be a totally non-meagre space which is connected and locally connected. Show that X cannot be the union of an infinite sequence (F \(_{n}\) ) of non-empty mutually disjoint closed sets. (Let H be the union of the frontiers H \(_{n}\) of the F \(_{n}\) in X; show that H is closed in X and that each set H \(_{n}\) is nowhere dense in H; to establish the latter point, consider a fundamental system of connected neighbourhoods of a point of \(\mathbf{H}_n\) , and argue by contradiction, using Chapter I, § 11, Proposition 3.)

\begin{enumerate}
\def\labelenumi{\arabic{enumi})}
\setcounter{enumi}{13}
\tightlist
\item
  Let X be the subspace of \(R^{2}\) consisting of the points \((r,0)\) , where r runs through the set Q of rational numbers, and the points \((k/n,1/n)\) , where n runs through the set of integers \(\geqslant 1\) and k runs through the set of all rational integers. Show that X is a Baire space but is not totally non-meagre.
\end{enumerate}

¶ 15) Let X be the subset of the plane \(\mathbf{R}^2\) consisting of the line \(\mathrm{D} = \{0\} \times \mathbf{R}\) and the points \((1/n, k/n^2)\) , where \(n\) runs through the set of integers \(>0\) and \(k\) runs through the set of all rational integers. a) For each point \((0, y)\) of D and each integer \(n > 0\) , let \(\mathrm{T}_n(y)\) be the set of points \((u, v) \in \mathbf{X}\) such that \(u \leqslant 1/n\) and \(|v - y| \leqslant u\) . Show that if we take as a fundamental system of neighbourhoods of each point \((0, y)\) the set of sets \(\mathrm{T}_n(y)\) , and as a fundamental system of neighbourhoods of each other point of X the set consisting of this point alone, then we have defined a topology \(\mathcal{G}\) on X, which is finer than the topology induced by that of \(\mathbf{R}^2\) , and with respect to which X is locally compact. b) Show that X is the union of a countable family of metrizable closed subspaces and that every point of X has a metrizable compact neighbourhood, but that X is not normal. {[}Let A be the set of all points \((0, y)\) where \(y\) is rational, and let \(\mathbf{B} = \mathbf{D} - \mathbf{A}\) ; show that every neighbourhood U of A in X meets every neighbourhood V of B in X; for this, consider for each \(n\) the set \(\mathbf{B}_n\) of points \(y \in \mathbf{B}\) such that \(\mathrm{T}_n(y) \subset \mathbf{V}\) , and note that the set B is not meagre in R.{]}

¶ 16) Let \(\mathbf{R}_{-}\) denote the topological space obtained by endowing the linearly ordered set \(\mathbf{R}\) of real numbers with the topology \(\mathcal{G}_{-}(\mathbf{R})\) (Chapter I, § 2, Exercise 5).

\begin{enumerate}
\def\labelenumi{\alph{enumi})}
\item
  Show that every open set in \(\mathbf{R}_{-}\) is the union of a countable family of mutually disjoint intervals which are either open or half-open on the left. Deduce that \(\mathbf{R}_{-}\) is perfectly normal (§ 4, Exercise 8) and totally non-meagre (Exercise 12).
\item
  Show that the space \(\mathbf{R}_{-} \times \mathbf{R}_{-}\) is not normal, and hence that \(\mathbf{R}_{-}\) is not metrizable, although \(\mathbf{R}_{-}\) contains a countable dense subset. (Considering the set of points \((x, y)\) in \(\mathbf{R}_{-} \times \mathbf{R}_{-}\) such that \(x + y = 1\) , show that \(\mathbf{R}_{-} \times \mathbf{R}_{-}\) has a closed subspace which is homeomorphic to the space \(\mathbf{X}\) defined in Exercise 15.)
\item
  Show that \(\mathbf{R}_{-}\) is a Lindelöf space (Chapter I, § 9, Exercise 15) and therefore is paracompact (§ 4, Exercise 23). (For each \(x \in \mathbf{R}_{-}\) , let \(\mathbf{U}_x\) be a semi-open interval {]} \(y(x)\) , \(x]\) with \(x\) as right-hand end-point. Show that the set of \(z \in \mathbf{R}_{-}\) which do not belong to any open interval {]} \(y(x)\) , \(x[\) is countable; use Chapter IV, § 2, Exercise 1.)
\item
  Show that every compact subset of \(R_{-}\) is countable {[}if A is a compact set in \(R_{-}\) , the topologies induced on A by \(\mathcal{G}_{-}(\mathbf{R})\) and \(\mathcal{G}_{0}(\mathbf{R})\) are identical; deduce that every point of A is the left-hand end-point of an interval contiguous to A{]}.
\end{enumerate}

\begin{enumerate}
\def\labelenumi{\arabic{enumi})}
\setcounter{enumi}{16}
\tightlist
\item
  \begin{enumerate}
  \def\labelenumii{\alph{enumii})}
  \tightlist
  \item
    Show that every product \(\mathbf{X} = \prod_{i\in \mathbf{I}}\mathbf{X}_i\) of complete metric spaces is a Baire space. {[}Argue as in the proof of Theorem 1, taking \(\mathbf{G}_n\) to be an elementary set (Chapter I, § 4, no. 1) all of whose projections are either the whole space or have a diameter \(\leqslant 1 / n\) ; then note that there is a countable subset J of I such that each \(\mathbf{G}_n\) is of the form \(\mathbf{H}_n\times \prod_{i\notin \mathbf{J}}\mathbf{X}_i\) , where \(\mathbf{H}_n\subset \prod_{i\in \mathbf{J}}\mathbf{X}_i.\) {]}
  \end{enumerate}
\end{enumerate}

\begin{enumerate}
\def\labelenumi{\alph{enumi})}
\setcounter{enumi}{1}
\tightlist
\item
  Show by an example that if I is uncountable then X need not be totally non-meagre (Exercise 12). {[}Use Exercises 1 and 23 b) of § 2.{]}
\end{enumerate}

¶ 18) a) Show that in a complete metric space, every perfect set contains a subset homeomorphic to Cantor's triadic set (Chapter IV, § 2, no. 5) (use Exercise 11 of Chapter IV, § 8).

\begin{enumerate}
\def\labelenumi{\alph{enumi})}
\setcounter{enumi}{1}
\item
  Deduce that, in a complete metric space with a countable base, every perfect set has the power of the continuum and that every closed set and every open set is either countable or has the power of the continuum {[}use Exercise 11 b){]}.
\item
  Show that in an uncountable complete metric space with a countable base, the set of perfect subsets has the power of the continuum. {[}Prove the assertion first for the compact space \(\{0, 1\}^{N}\) , then use a) and § 2, Exercise 12 b).{]}
\item
  Let X be an uncountable complete metric space with a countable base. Show that there is a subset Z of X such that both Z and X---Z have the power of the continuum and such that neither Z nor X---Z contains a non-empty perfect set. {[}Use b) and c) and the method described in Set Theory, Chapter III, § 6, Exercise 24 a).{]} If, furthermore, X has no isolated point, show that neither Z nor X---Z is meagre. (If A is a countable intersection of dense open subsets of X, show that A contains non-empty perfect sets by using c) and § 2, Exercise 24.) Deduce that in this case Z is not an almost open set in X (Exercise 6; note that the reasoning above shows that if U is any non-empty open subset of X, U ∩ Z is not meagre).
\end{enumerate}

\begin{enumerate}
\def\labelenumi{\arabic{enumi})}
\setcounter{enumi}{18}
\tightlist
\item
  Let A be a dense subspace of a topological space X, and let f be a continuous mapping of A into a complete metric space Y. Show that the set of points of X where f has no limit (relative to A) is a meagre set in X {[}consider, for each n \textgreater{} 0, the set of points \(x \in E\) at which the oscillation \(\omega(x; f)\) of f is \textless{} 1/n{]}.
\end{enumerate}

¶ 20) Let \(f\) be a homeomorphism of a subset A of a complete metric space X onto a subset \(\mathbf{A}'\) of a complete metric space \(\mathbf{X}'\) . Show that there exists an extension \(\overline{f}\) of \(f\) to a set B which is a countable intersection of open sets in X, such that \(f\) is a homeomorphism of B onto a set \(\mathbf{B}'\) which is a countable intersection of open sets in \(\mathbf{X}'\) (apply Exercise 19 to \(f\) and to the inverse homeomorphism \(g\) ).

¶ 21) a) Let \((f_n)\) be a sequence of continuous mappings of a topological space X into a perfectly normal space Y (§ 4, Exercise 8) such that, for each \(x \in \mathbf{X}\) , the sequence \((f_n(x))\) has a limit \(f(x)\) in Y. Show that, for every closed subset S of Y, \(A = \overline{f}(S)\) is a countable intersection of open sets. {[}There exists a decreasing sequence \((U_n)\) of open sets in Y such that \(S = \bigcap_{n} \overline{U}_n\) ; note that \(f(x) \in S\) if and only if, for each integer \(n > 0\) , there exists an integer \(k \geqslant 0\) such that \(f_{n+k}(x) \in U_n\) .{]} Deduce that \(\overline{A} - A\) is meagre in X.

\begin{enumerate}
\def\labelenumi{\alph{enumi})}
\setcounter{enumi}{1}
\tightlist
\item
  Deduce from a) that if we suppose in addition that Y is a metrizable space of countable type, then the set of all \(x \in X\) such that f is not continuous at x is meagre. {[}If \((\mathbf{V}_{n})\) is a base of the topology of Y and if \(S_{n} = Y - V_{n}\) and \(\mathbf{A}_{n} = f^{-1}(\mathbf{S}_{n})\) , show that x must belong to one of the sets \(\overline{\mathbf{A}}_{n} - \mathbf{A}_{n}\) .{]}
\end{enumerate}

\begin{enumerate}
\def\labelenumi{\arabic{enumi})}
\setcounter{enumi}{21}
\tightlist
\item
  \begin{enumerate}
  \def\labelenumii{\alph{enumii})}
  \tightlist
  \item
    Let X be a Baire space, Y a metric space, \((f_{n})\) a sequence of continuous mappings of X into Y such that, for each \(x \in X\) , the sequence \((f_{n}(x))\) has a limit \(f(x)\) in Y. Then the set of all points \(x \in X\) at which f is not continuous is meagre. {[}Prove that, for each integer n \textgreater{} 0, the set of all points \(x \in X\) at which the oscillation \(\omega(x; f)\) of f is \(\leqslant 1/n\) contains a dense open set; for this, if \(G_{p}\) denotes the set of \(x \in X\) for which the distance from \(f_{p}(x)\) to \(f_{q}(x)\) is \(\leqslant 1/2n\) for all \(q \geqslant p\) , show that the union of the interiors of the sets \(G_{p}\) is a dense open set.{]}
  \end{enumerate}
\end{enumerate}

\begin{enumerate}
\def\labelenumi{\alph{enumi})}
\setcounter{enumi}{1}
\tightlist
\item
  Let K be the compact interval {[}0, 1{]} in R. Give an example of a sequence of continuous mappings \(f_{n}\) of K into the compact space \(K^{K}\) such that \(\lim_{n\to\infty}f_{n}(x)=f(x)\) exists for each \(x\in K\) , but such that f is not continuous at any point of K.
\end{enumerate}

¶ 23) Let X be a topological space, and let Y, Z be two metric spaces, with metrics d, d' respectively. Let f: X × Y → Z be a mapping such that y → f(x₀, y) is continuous on Y for each x₀ ∈ X and such that x → f(x, y₀) is continuous on X for each y₀ ∈ Y.

\begin{enumerate}
\def\labelenumi{\alph{enumi})}
\item
  For each real number \(\varepsilon > 0\) , each \(b \in \mathbf{Y}\) and each \(x \in \mathbf{X}\) , let \(g(x; b, \varepsilon)\) denote the least upper bound of the numbers \(\alpha > 0\) such that the relation \(d(b, y) < \alpha\) implies \(d'(f(x, b), f(x, y)) \leqslant \varepsilon\) . Show that the function \(x \to g(x; b, \varepsilon)\) is upper semi-continuous.
\item
  If X is a Baire space, deduce from a) and Theorem 2 that for each \(b \in Y\) there is a set \(S_{b}\) in X, whose complement is meagre, such that f is continuous at \((a, b)\) for each \(a \in S_{b}\) .
\end{enumerate}

\begin{enumerate}
\def\labelenumi{\arabic{enumi})}
\setcounter{enumi}{23}
\tightlist
\item
  Let X, Y be two topological spaces. A mapping \(f: \mathbf{X} \to \mathbf{Y}\) is said to be almost open if, for every closed subset S of Y, \(\overline{f}(\mathbf{S})\) is almost open in X (Exercise 6).
\end{enumerate}

\begin{enumerate}
\def\labelenumi{\alph{enumi})}
\item
  Let \(g: \mathbf{X} \to \mathbf{Y}\) be a mapping with the property that there exists a meagre set \(\mathbf{M}\) in \(\mathbf{X}\) such that \(g|_{\mathbf{CM}}\) is continuous. Show that \(g\) is almost open. Consider the converse, when the topology of \(\mathbf{Y}\) has a countable base.
\item
  Suppose that Y is perfectly normal (§ 4, Exercise 8). Let \((f_n)\) be a sequence of almost open mappings of X into Y such that
\end{enumerate}

\[
\lim _ {n \rightarrow \infty} f _ {n} (x) = f (x)
\]

exists for each \(x \in \mathbf{X}\) . Show that \(f\) is almost open {[}argue as in Exercise 21 a){]}.

\begin{enumerate}
\def\labelenumi{\arabic{enumi})}
\setcounter{enumi}{24}
\tightlist
\item
  Let R be an open equivalence relation on a topological space X. Show that if X is non-meagre (Exercise 7) {[}resp. a Baire space, resp. totally non-meagre (Exercise 12){]}, then so is X/R.
\end{enumerate}

¶ 26) Let X be a locally compact metrizable space, let d be a metric compatible with the topology of X, and let R be a closed equivalence relation on X, all of whose equivalence classes are compact. Then X/R is locally compact (Chapter I, § 10, Proposition 9) and metrizable (§ 4, Exercise 24).

\begin{enumerate}
\def\labelenumi{\alph{enumi})}
\tightlist
\item
  For each pair of points \(u, v\) of \(\mathbf{X} / \mathbf{R}\) , set
\end{enumerate}

\[
\rho (u, v) = \sup _ {\varphi (x) = u} d (x, \overline {{\varphi}} ^ {1} (v)),
\]

where \(\varphi : \mathbf{X} \to \mathbf{X} / \mathbf{R}\) is the canonical mapping, and for each \(u \in \mathbf{X} / \mathbf{R}\) let

\[
\lambda (u) = \lim _ {v \rightarrow u, v \neq u} \sup _ {\rho (u, v)}.
\]

Show that \(\lambda\) is upper semi-continuous on X/R and that, for each \(\alpha > 0\) , the set of points \(u \in \mathrm{X} / \mathrm{R}\) such that \(\lambda(u) \geqslant \alpha\) is nowhere dense. {[}Show that otherwise there would exist a compact subset K of X and an infinite sequence \((x_n)\) of points of K such that \(d(x_m, x_n) > \frac{1}{2} \alpha\) for \(m \neq n\) .{]} b) Deduce from a) that there is a dense subset M of X/R which is a countable intersection of open sets in X/R and is such that for each \(x \in \varphi^{-1}(M)\) the image under \(\varphi\) of every neighbourhood of x in X is a neighbourhood of \(\varphi(x)\) in X/R (use Theorem 1).

\begin{enumerate}
\def\labelenumi{\arabic{enumi})}
\setcounter{enumi}{26}
\tightlist
\item
  \begin{enumerate}
  \def\labelenumii{\alph{enumii})}
  \tightlist
  \item
    Let G be a topological group, and let A be an almost open set (Exercise 6) in G. Show that if A is not meagre then AA \(^{-1}\) is a neighbourhood of the identity element in G. {[}Note first that G is a Baire space by virtue of Exercise 3. Let A* be the union of all open sets U ⊂ G such that U ∩ C A is meagre. Show that A* is not empty and that if x ∈ G is such that xA* meets A*, then xA meets A. Hence show that AA \(^{-1}\) ⊃ A*(A*) \(^{-1}\) .{]}
  \end{enumerate}
\end{enumerate}

\begin{enumerate}
\def\labelenumi{\alph{enumi})}
\setcounter{enumi}{1}
\item
  Deduce from a) that an almost open subgroup of a topological group G is either meagre or else open (and closed). In particular, a countable topological group which is a Hausdorff Baire space is discrete.
\item
  Show that there exist subgroups H of the topological group R such that R/H is countably infinite (use a Hamel base); such a subgroup is dense but not meagre nor almost open.
\item
  Let G be a metrizable topological group whose left and right uniformities coincide. Show that if there exists a metric on G, compatible with the topology of G, and with respect to which G is a complete metric space, then G is a complete metrizable group {[}use b); § 3, Exercise 2; and § 2, Exercise 24{]}.
\end{enumerate}

\begin{enumerate}
\def\labelenumi{\arabic{enumi})}
\setcounter{enumi}{27}
\tightlist
\item
  Let \(\mathbf{G}, \mathbf{G}'\) be two topological groups. Show that every continuous homomorphism \(f\) of \(\mathbf{G}\) onto \(\mathbf{G}'\) is a strict morphism in each of the following two cases:
\end{enumerate}

\begin{enumerate}
\def\labelenumi{\alph{enumi})}
\item
  G is locally compact and \(\sigma\) -compact, and \(\mathbf{G}'\) is non-meagre {[}show that there is a relatively compact open set U in G such that the interior of \(f(\overline{\mathbf{U}})\) is not empty; deduce that for each compact neighbourhood V of the identity element of G, the interior of \(f(\mathbf{V})\) is not empty, and then apply Chapter III, § 2, Exercise 18{]}.
\item
  G is complete, the topology of G has a countable base, and \(G'\) is non-meagre. {[}By using the fact that, for each neighbourhood U of the identity element e of G, G is the union of a sequence of sets of the form \(x_{n}U\) , show that for each open set A in G there is an open set \(A'\) in \(G'\) which contains \(f(A)\) and is such that \(f(A)\) is dense in \(A'\) . Then consider a fundamental system \((\mathbf{U}_{n})\) of symmetric neighbourhoods of e in G, such that \(\overline{U}_{n+1}^{2} \subset U_{n}\) , and show that \(f(\mathbf{U}_{p})\) contains the open set \(U_{p+1}'\) ; to do this, take any point \(a' \in U_{p+1}'\) and then construct by induction a sequence \((b_{n})\) of points of G such that \(b_{n} \in U_{p+n}\) and such that \(f(b_{1}b_{2}\ldots b_{n})\) tends to \(a'\) ; note that, if \(x' \in U_{k}'\) , then there exists \(y \in \mathbf{U}_k\) such that \(f(y) \in x' \mathrm{U}_{k+1}'\) .{]} The hypothesis that \(\mathbf{G}\) has a countable base is essential, as is shown by the example where \(\mathbf{G}'\) is the group \(\mathbf{R}\) with the topology of the real line, \(\mathbf{G}\) is the group \(\mathbf{R}\) with the discrete topology and \(f: \mathbf{G} \to \mathbf{G}'\) is the identity mapping.
\end{enumerate}

\begin{enumerate}
\def\labelenumi{\arabic{enumi})}
\setcounter{enumi}{28}
\tightlist
\item
  Let G be a locally compact, σ-compact, topological group, or else a complete metrizable group whose topology has a countable base. Let H be a closed normal subgroup of G and let A be a closed subgroup of G. Then the quotient group A/(A∩H) is isomorphic to AH/H if and only if AH is a closed subgroup of G {[}to show that the condition is necessary, note that A/(A∩H) is complete, by Chapter III, § 4, no. 6, Proposition 13, and Chapter IX, § 3, no. 1, Proposition 4; to show that the condition is sufficient, use Exercise 28{]}.
\end{enumerate}

¶ 30) Let G be a group and let d be a metric on G such that, with respect to the topology G defined by d on G, the mappings \(y \to x_{0}y\) and \(y \to yx_{0}\) of G into G are continuous for each \(x_{0} \in G\) .

\begin{enumerate}
\def\labelenumi{\alph{enumi})}
\item
  Show that if \(\mathbf{G}\) is complete with respect to the topology \(\mathfrak{C}\) , then the mapping \((x,y)\to xy\) is continuous on \(\mathbf{G}\times\mathbf{G}\) (use Exercise 23).
\item
  Suppose in addition that \(\mathcal{C}\) has a countable base. Show that the mapping \(x \to x^{-1}\) is then continuous on \(\mathbf{G}\) , i.e., that the topology \(\mathcal{C}\) is compatible with the group structure of \(\mathbf{G}\) . {[}In the group \(\mathbf{G} \times \mathbf{G}\) , endowed with the topology \(\mathcal{C} \times \mathcal{C}\) , consider the set \(\mathbf{F}\) of all points \((x, x^{-1})\) where \(x \in \mathbf{G}\) ; \(\mathbf{F}\) is closed in \(\mathbf{G} \times \mathbf{G}\) , and the law of composition \((x, x^{-1})(y, y^{-1}) = (xy, y^{-1}x^{-1})\) defines a group structure on \(\mathbf{F}\) ; show that the topology induced on \(\mathbf{F}\) by that of \(\mathbf{G} \times \mathbf{G}\) is compatible with this group structure by using \(a\) ). Then prove, by arguing as in Exercise 28 \(b\) , that the projection \(\mathrm{pr}_1\) of \(\mathbf{F}\) onto \(\mathbf{G}\) is bicontinuous.{]}
\end{enumerate}

